\documentclass[12pt, uplatex, dvipdfmx, article]{jlreq}
\usepackage{amssymb, amsmath, latexsym}
\usepackage[skins]{tcolorbox}
% \usepackage[haranoaji]{pxchfon} % 原ノ味フォントを使う


% 5つの引数を受け取る \present コマンドの定義
\newcommand{\speak}[5]{ %内部的にtcolorboxを使っている
\begin{tcolorbox}[
    colback= white,        %背景色を白
    colframe= white,       %枠線を白色に
    colbacktitle= white,
    before skip= 1pt,      %ボックス外側の「上」余白
    after skip= 1pt,       %ボックス外側の「下」余白
    top= 0.7mm,            %ボックス内側の「上」余白
    bottom= 0.7mm,         %ボックス内側の「下」余白
    sidebyside,            %ボックスを左右に分割
    sidebyside align= top, %上端を揃える
    lefthand width= 28mm,  %左側（時間）の幅
    sidebyside gap= 3mm,   %時間と講演者の間の隙間
    lower separated= false %中央の区切り線を消す
    ]
    {#1}\hspace{1.5mm}～\hspace{1.5mm}{#2} %#1：開始時間、#2：終了時間
    \tcblower                 % ここから右側（講演者・タイトル）
    {#3}(#4)\\ %#3：講演者名、#4：所属・指導教員名学年・分野
    {#5} %5：講演タイトル
\end{tcolorbox}
}


\begin{document}
\pagestyle{empty} %ページ番号を消す

\begin{center}
{\huge %研究集会名
    Workshop on Mathematical Research\&Study Conference in Hakkaisan 2027
}
\vspace{12pt}

%他詳細
\begin{tabular}{l}
    {日程：2027年1月3日 12:00～18:00}\\
    {会場：日本大学　八海山セミナーハウス}\\
    {形式：対面/オンラインハイブリッド開催(オンラインは聴講のみ)}\\
    {Zoomリンク：}
\end{tabular}

\vspace{10pt}
    プログラム
\end{center}

% \speak{開始時間}{終了時間}{講演者}{所属}{学年}{講演タイトル}


% \speak{12:00}{13:00}{山田太郎}{A大学}{M1}{微分方程式について}
% \speak{13:00}{14:00}{佐藤花子}{B大学}{B4}{代数方程式}

\speak{12:00}{12:20}{川村 聡一郎}{日大理工院・水野研M1・応用}{Allen-Cahn方程式の導出とパラメータの関係付け}

\speak{12:30}{12:50}{田中 涼太}{日大理工院・月岡研M1・代数}{いろんなイデアルの対応関係}

\speak{13:00}{13:20}{川村 聡一郎}{日大理工院・水野研M1・応用}{最近の数学界隈を騒がせる定理証明支援系-Lean-とは？}

\speak{13:30}{13:50}{宮崎 裕貴}{日大理工院・伊藤研M2・応用}{TBA}


%講演未定者 院生
\speak{hh:mm}{hh:mm}{鈴木 輝星}{日大理工院・平石研M1・応用}{TBA}
\speak{hh:mm}{hh:mm}{岡田 晟}{日大理工院・水野研M1・応用}{TBA}
\speak{hh:mm}{hh:mm}{鈴木 聖也}{中大XX理工院・XXM1・解析}{TBA}

\speak{hh:mm}{hh:mm}{HU BOQING}{日大理工院・青柳研M1・応用}{TBA}
\speak{hh:mm}{hh:mm}{江原 響}{日大理工院・小紫研M1・応用}{TBA}
\speak{hh:mm}{hh:mm}{廣瀬 友大}{日大理工院・小紫研M1・応用}{TBA}
\speak{hh:mm}{hh:mm}{小久保 勇亮}{日大理工院・伊藤研M1・応用}{TBA}
\speak{hh:mm}{hh:mm}{相馬 悠人}{日大理工院・石原研M1・応用}{TBA}
\speak{hh:mm}{hh:mm}{松山 由佳}{日大理工院・青柳研M1・応用}{TBA}
\speak{hh:mm}{hh:mm}{三瓶 那々花}{日大理工院・石原研M1・応用}{TBA}
\speak{hh:mm}{hh:mm}{斑目 穂高}{日大理工院・小紫研M1・応用}{TBA}
\speak{hh:mm}{hh:mm}{小野 雅弘}{日大理工院・善本研M1・応用}{TBA}
\speak{hh:mm}{hh:mm}{中村 京志郎}{日大理工院・平石研M1・応用}{TBA}
\speak{hh:mm}{hh:mm}{丸山 晃汰}{日大理工院・平田研M1・応用}{TBA}
\speak{hh:mm}{hh:mm}{菊池 健太}{日大理工院・橋口研M1・応用}{TBA}


%講演未定者 学部生
\speak{hh:mm}{hh:mm}{栗田}{日大理工・平石研B4・応用}{TBA}
\speak{hh:mm}{hh:mm}{末次}{日大理工・青柳研B4・応用}{TBA}
\speak{hh:mm}{hh:mm}{岩橋 優里伽}{日大理工・水野研B4・解析}{TBA}
\speak{hh:mm}{hh:mm}{心春}{日大理工・小紫研B4・応用}{TBA}




%これより上に講演内容を書く

\begin{flushleft} %世話人の名前
    世話人：川村聡一郎(日大理工院・M1)
\end{flushleft}
\end{document}
